\section{High parameter derivatives of Stieltjes constants}
\label{sec:asymptotics}

The parameter-derivative tower admits a useful separation into a
finite polynomial coefficient and an exponentially small spectral
tail. This gives a different asymptotic regime from the frequently
studied limit of a large Stieltjes index with derivative order fixed.
Here the derivative order \(k\) tends to infinity, while the index
\(n\) may grow proportionally to \(\log k\). The resulting profile is
the reciprocal gamma function.

\subsection{An exact spectral expansion}

Fix the Laurent convention
\begin{equation}\label{eq:stieltjes-convention}
\zeta(s,a)=\frac{1}{s-1}
 +\sum_{n=0}^{\infty}\frac{(-1)^n}{n!}\gamma_n(a)(s-1)^n,
\qquad a>0.
\end{equation}
The superscript \(\gamma_n^{(k)}(a)\) always denotes differentiation
with respect to \(a\), not to the index or the zeta variable. Define
\begin{equation}\label{eq:coefficient-polynomial}
P_{n,k}(x)=n![z^n]\left(e^{xz}\prod_{j=1}^{k}
                  \left(1+\frac zj\right)\right).
\end{equation}
This is a degree-\(n\) polynomial in \(x\); its coefficients are
rational. In particular \(P_{0,k}=1\) and
\(P_{1,k}(x)=x+H_k\).

\begin{theorem}[Spectral expansion and a finite-tail certificate]
\label{thm:spectral}
For \(a>0\), integers \(k\ge1\), and \(n\ge0\),
\begin{equation}\label{eq:spectral}
\frac{(-1)^{n+k}}{k!}\gamma_n^{(k)}(a)
=\sum_{m=0}^{\infty}
 \frac{P_{n,k}(-\log(a+m))}{(a+m)^{k+1}}.
\end{equation}
The series converges absolutely, locally uniformly for \(a>0\).
If the sum is truncated before \(m=M\), where \(b=a+M\ge1\),
then for every real \(0<R<k\) the magnitude of the omitted part is at
most
\begin{equation}\label{eq:spectral-tail}
\frac{n!}{R^n}\prod_{j=1}^{k}\left(1+\frac Rj\right)
\left(b^{-k-1+R}+\frac{b^{-k+R}}{k-R}\right).
\end{equation}
\end{theorem}

\begin{proof}
The Hurwitz-zeta parameter identity
\[
\partial_a^k\zeta(s,a)=(-1)^k(s)_k\zeta(s+k,a)
\]
follows initially by differentiating its absolutely convergent
Dirichlet series and then by analytic continuation; here
\((s)_k=s(s+1)\cdots(s+k-1)\).
The pole in \eqref{eq:stieltjes-convention} is independent of \(a\).
Put \(s=1+z\), use
\((1+z)_k=k!\prod_{j=1}^k(1+z/j)\), and compare coefficients.
For \(|z|<k\), the series for \(\zeta(k+1+z,a)\) converges normally.
This gives \eqref{eq:spectral} and justifies the coefficient
interchange. Alternatively, absolute convergence follows directly
because each summand is a linear combination of
\((a+m)^{-k-1}\log^\ell(a+m)\), \(0\le\ell\le n\).
The parameter identity is also the starting point of
Coffey's all-order formulas~\cite{Coffey}.

On the circle \(|z|=R\), and for \(x\ge b\ge1\),
\[
\left|e^{-z\log x}\prod_{j=1}^k(1+z/j)\right|
\le x^R\prod_{j=1}^k(1+R/j).
\]
Cauchy's coefficient estimate bounds the \(m\)-th omitted summand
by \(n!R^{-n}\prod_j(1+R/j)(a+m)^{-k-1+R}\).
For the positive decreasing power, comparison with its integral gives
\[
\sum_{m=M}^{\infty}(a+m)^{-k-1+R}
\le b^{-k-1+R}+\int_b^\infty x^{-k-1+R}\,dx.
\]
This is \eqref{eq:spectral-tail}.
\end{proof}

The condition \(b\ge1\) belongs to the stated majorant, rather than to
the exact expansion. It always holds for the tail after the first
term, because then \(b=a+1>1\). The choice of \(R\) can be optimized
for a given \(n,k,M\); there is no need to differentiate a
highly oscillatory numerical representation of \(\gamma_n\).

\subsection{Exact rational enclosures at parameter one}

Write \(\stirling{v}{u}\) for the unsigned Stirling number of the
first kind. The elementary identity
\begin{equation}\label{eq:stirling-product}
\prod_{j=1}^k(1+z/j)
 =\frac1{k!}\sum_{n=0}^k\stirling{k+1}{n+1}z^n
\end{equation}
turns the first spectral term at \(a=1\) into an exact rational
number.

\begin{corollary}[Rational certificates]\label{cor:rational}
Let \(k\ge2\), \(n\ge0\), and
\[
e_n(k)=
\begin{cases}
\stirling{k+1}{n+1}/k!,&0\le n\le k,\\
0,&n>k.
\end{cases}
\]
For any integer \(1\le h<k\), set
\begin{equation}\label{eq:rational-radius}
B_{k,n,h}
=h^{-n}\binom{k+h}{h}
 \left(2^{-k-1+h}+\frac{2^{-k+h}}{k-h}\right).
\end{equation}
Then the following is an exact inclusion between real numbers:
\begin{equation}\label{eq:rational-inclusion}
\frac{(-1)^{n+k}\gamma_n^{(k)}(1)}{k!\,n!}
\in[e_n(k)-B_{k,n,h},\,e_n(k)+B_{k,n,h}].
\end{equation}
One may take the minimum radius over all \(h=1,\ldots,k-1\).
In particular, the simple choice \(h=1\) gives
\begin{equation}\label{eq:simple-rational-radius}
B_k=(k+1)\left(2^{-k}+\frac{2^{-k+1}}{k-1}\right),
\end{equation}
independently of \(n\).
\end{corollary}

\begin{proof}
Apply Theorem~\ref{thm:spectral} with \(a=M=1\), \(b=2\), and
\(R=h\), and divide its bound by \(n!\). The product in that bound is
\(\prod_{j=1}^k(1+h/j)=\binom{k+h}{h}\).
Equation~\eqref{eq:stirling-product} supplies the center.
All quantities in the resulting endpoints are rational.
\end{proof}

Thus the term \emph{certificate} has its literal meaning here:
integer arithmetic constructs the endpoints, and the proved tail
bound supplies the inclusion. Floating-point agreement is not being
renamed a proof. For example, positivity of an endpoint
\(e_n(k)-B_{k,n,h}\) rigorously proves
\((-1)^{n+k}\gamma_n^{(k)}(1)>0\).
The accompanying script emits centers, radii, and both endpoints
as numerator--denominator pairs, and checks their arithmetic exactly.

\subsection{A uniform coefficient lemma}

The main asymptotic step is an elementary coefficient estimate.
It is convenient to set
\(\fall{n}{j}=n(n-1)\cdots(n-j+1)\) for \(j\le n\), and
\(\fall{n}{j}=0\) for \(j>n\).

\begin{lemma}[Two uniform coefficient terms]\label{lem:coefficient}
Suppose \(f\) is analytic on a neighborhood of the closed disk
\(|z|\le R\), with \(|f(z)|\le M\) there. Let \(L>0\),
\(n\ge1\), and \(r=n/L<R\). Put
\[
Q_n(L;f)=\frac{n!}{L^n}[z^n]e^{Lz}f(z).
\]
Then
\begin{align}
|Q_n(L;f)-f(r)|
&\le\frac{Mn}{L^2R^2(1-r/R)^3},\label{eq:coefficient-first}\\
\left|Q_n(L;f)-f(r)+\frac{n}{2L^2}f''(r)\right|
&\le\frac{3Mn}{L^3R^3(1-r/R)^5}.\label{eq:coefficient-second}
\end{align}
For \(n=0\), \(Q_0(L;f)=f(0)\) exactly.
\end{lemma}

\begin{proof}
Write \(f(z)=\sum_{j\ge0}c_jz^j\). Since \(|c_j|\le MR^{-j}\),
\[
Q_n(L;f)=\sum_{j\ge0}c_j\frac{\fall{n}{j}}{L^j}
\]
is a finite sum, which we compare with the absolutely convergent
Taylor series at \(r\).

For completeness, the needed inequalities remain valid when \(j>n\).
Choose \(j\) independent uniform elements of a set of \(n\) elements.
The probability of no repeated choice is
\(\fall{n}{j}/n^j\), including its value zero when \(j>n\).
For the \(K=\binom j2\) pair-collision events, every event has
probability \(1/n\), and every intersection of two distinct events
has probability \(1/n^2\), even if the pairs share a position.
The union bound and the first two Bonferroni inequalities give
\begin{align*}
0&\le n^j-\fall{n}{j}\le\binom j2 n^{j-1},\\
0&\le\fall{n}{j}-n^j+\binom j2 n^{j-1}
 \le\binom{\binom j2}{2}n^{j-2}.
\end{align*}
At \(j=0,1,2\), the second remainder is zero, so no negative-power
exception is needed. Now use
\[
\sum_{j\ge2}\binom j2 x^{j-2}=(1-x)^{-3},
\qquad
\sum_{j\ge3}\binom{\binom j2}{2}x^{j-3}=3(1-x)^{-5}.
\]
The first inequality gives \eqref{eq:coefficient-first}.
In the second comparison,
\[
\sum_{j\ge2}c_j\binom j2\frac{n^{j-1}}{L^j}
=\frac{n}{2L^2}f''(r).
\]
This proves \eqref{eq:coefficient-second}.
\end{proof}

\subsection{The reciprocal-gamma profile}

\begin{theorem}[A joint logarithmic-range asymptotic]\label{thm:profile}
Fix \(0<a_0\le a_1<\infty\) and \(C>0\). For
\(a\in[a_0,a_1]\), put
\[
L=\log(k/a),\qquad r=\frac nL,\qquad
g(r)=\frac1{\Gamma(1+r)},\qquad
\rho=\frac{a_1}{a_1+1}<1.
\]
Uniformly over integers \(0\le n\le CL\), as \(k\to\infty\),
\begin{equation}\label{eq:uniform-profile}
\frac{(-1)^{n+k}a^{k+1}\gamma_n^{(k)}(a)}{k!\,L^n}
=g(r)-\frac{n}{2L^2}g''(r)
 +O_{a_0,a_1,C}\left(\frac n{L^3}+k^{C+1}\rho^k\right).
\end{equation}
Equivalently, the two displayed terms are
\begin{equation}\label{eq:profile-polygamma}
\frac1{\Gamma(1+r)}
\left[1+\frac{n}{2L^2}
 \left\{\psi^{(1)}(1+r)-\psi(1+r)^2\right\}\right].
\end{equation}
The formula includes \(n=0\), for which the first spectral term is
exactly one. In particular, if \(n/\log(k/a)\to c\ge0\) in a bounded
range, then the left side of \eqref{eq:uniform-profile} tends to
\(1/\Gamma(1+c)\).
\end{theorem}

\begin{proof}
Separate the \(m=0\) term of \eqref{eq:spectral}. Its generating
function factors as
\[
e^{-z\log a}\prod_{j=1}^k(1+z/j)=e^{Lz}f_k(z),
\qquad
f_k(z)=k^{-z}\prod_{j=1}^k(1+z/j).
\]
The gamma product gives
\[
f_k(z)=\frac{\Gamma(k+1+z)}
 {\Gamma(k+1)k^z\Gamma(1+z)}
      =g(z)e^{D_k(z)}.
\]
On \(|z|<k+1\), define the analytic function
\begin{equation}\label{eq:Dk}
D_k(z)=z(H_k-\log k-\gamma)
 +\sum_{\ell=2}^{\infty}
   \frac{(-1)^\ell z^\ell}{\ell}
       \bigl(\zeta(\ell)-H_k^{(\ell)}\bigr).
\end{equation}
The factorization follows first on a small disk from the
Weierstrass product and then throughout \(|z|<k+1\) by analytic
continuation. Defining \(D_k\) by its series avoids taking the
logarithm of \(g\) at its zeros.

The elementary bounds
\[
0<H_k-\log k-\gamma<\frac1{2k},\qquad
\sum_{j>k}j^{-2}\le\frac1k
\]
imply, for \(|z|\le R<k+1\),
\begin{equation}\label{eq:Dk-bound}
|D_k(z)|\le\frac{R}{2k}
   +\frac{R^2}{2k(1-R/(k+1))}.
\end{equation}
Indeed, for \(\ell\ge2\), bound the tail of \(\zeta(\ell)\) by
\((k+1)^{2-\ell}/k\), and use \(1/\ell\le1/2\).
Consequently \(f_k\) is uniformly bounded on every fixed disk,
\(f_k-g=O_R(k^{-1})\), and Cauchy's estimate gives
\(f_k''-g''=O_C(k^{-1})\) on \([0,C]\).
The same estimate with \(R=r\), or the real product, also gives the
slightly sharper
\[
f_k(r)-g(r)=O_C(r/k)\qquad(0\le r\le C).
\]
It is this factor \(r\) that preserves the exact \(n=0\) first term.

Choose \(R=C+1\) in Lemma~\ref{lem:coefficient}. For \(n\ge1\)
the normalized first spectral term is
\[
Q_n(L;f_k)
=f_k(r)-\frac{n}{2L^2}f_k''(r)+O_C(n/L^3).
\]
Replacing \(f_k\) by \(g\) costs
\(O_C(n/(kL)+n/(kL^2))\), which is absorbed in \(O(n/L^3)\)
because \(L^2\le k\) for all sufficiently large \(k\), uniformly
on the compact parameter interval.

It remains to control the other spectral terms. Apply
\eqref{eq:spectral-tail} with \(M=1\), and multiply the bound by
\(a^{k+1}/L^n\). The result is
\[
\frac{n!}{(RL)^n}\prod_{j=1}^k(1+R/j)
\left(\frac{a}{a+1}\right)^{k+1}(a+1)^R
\left(1+\frac{a+1}{k-R}\right).
\]
The first factor is at most one: for \(n\ge1\), use
\(n!\le n^n\) and \(n/L\le C<R\); for \(n=0\), it is one by
convention. The product is \(O_R(k^R)\), as follows already from
\eqref{eq:Dk-bound} at \(z=R\).
The remaining factors are uniformly bounded apart from
\((a/(a+1))^{k+1}\le\rho^{k+1}\). Thus the normalized tail is
\(O_{a_0,a_1,C}(k^{C+1}\rho^k)\). The case \(n=0\) follows
directly from \(P_{0,k}=1\).
Finally,
\(g''(r)=g(r)\{\psi(1+r)^2-\psi^{(1)}(1+r)\}\).
\end{proof}

\begin{corollary}[A uniform eventual sign]\label{cor:sign}
For every compact interval \([a_0,a_1]\subset(0,\infty)\) and
every \(C>0\), there is \(k_0\) such that
\[
(-1)^{n+k}\gamma_n^{(k)}(a)>0
\]
for all \(k\ge k_0\), \(a\in[a_0,a_1]\), and integers
\(0\le n\le C\log(k/a)\).
\end{corollary}

\begin{proof}
The continuous function \(g\) is strictly positive on \([0,C]\),
so its minimum there is positive. Both the displayed correction
and the remainder in Theorem~\ref{thm:profile} tend uniformly to zero.
The other factors in the normalization are positive.
\end{proof}

The corollary is restricted to its stated range. It gives no sign
theorem when \(n\) grows arbitrarily quickly compared with \(\log k\).
For a concrete finite pair \(n,k\) at \(a=1\), the rational interval
in Corollary~\ref{cor:rational} is a separate, effective certificate.

\begin{figure}[tbp]
\centering
\includegraphics[width=\textwidth]{figures/derivative-profile.pdf}
\caption{The reciprocal-gamma profile at \(a=1\).
The first spectral term is computed from the elementary-symmetric
coefficient recurrence; its rigorously bounded spectral error is
smaller than the plotted resolution.
The lower panel compares its deviation from the limiting profile
with its deviation after the explicit second term is included.
These plotted arithmetic evaluations are floating point; the separate
rational certificates do not depend on them.}
\label{fig:profile}
\end{figure}

\subsection{Fixed-index expansions and the classical mechanism}

For a fixed \(n\), it is also useful to retain the polynomial in
\(L=\log(k/a)\). Define Appell polynomials
\[
A_n(L)=n![z^n]\frac{e^{Lz}}{\Gamma(1+z)}.
\]
The first few are
\begin{align*}
A_0(L)&=1,\\
A_1(L)&=L+\gamma,\\
A_2(L)&=(L+\gamma)^2-\zeta(2),\\
A_3(L)&=(L+\gamma)^3-3\zeta(2)(L+\gamma)+2\zeta(3).
\end{align*}
The standard gamma-ratio expansion~\cite{DLMF} gives, uniformly on
fixed \(z\)-disks,
\[
e^{D_k(z)}=
1+\frac{z+z^2}{2k}
 +\frac{3z^4+2z^3-3z^2-2z}{24k^2}+O(k^{-3}).
\]
Consequently, for fixed \(n\ge1\) and \(a\) in a compact positive
interval,
\begin{align}
\frac{(-1)^{n+k}a^{k+1}\gamma_n^{(k)}(a)}{k!}
={}&A_n(L)
 +\frac{nA_{n-1}(L)+n(n-1)A_{n-2}(L)}{2k}\notag\\
&+O_n\left(\frac{(1+L)^{n-1}}{k^2}\right)
 +O_{a_0,a_1,n}(k^2\rho^k).
\label{eq:fixed-index}
\end{align}
Terms with a negative Appell index are omitted.
To see the stated polynomial error, note that the analytic
remainder after the first correction vanishes at \(z=0\);
its coefficient convolution with \(e^{Lz}g(z)\) therefore has
degree at most \(n-1\) in \(L\). For \(n=0\) the first term is
exact and only the spectral error remains.
Higher powers of \(k^{-1}\) follow by the same coefficient procedure.
For example the \(k^{-2}\) numerator is
\[
\frac1{24}\left(
3\fall{n}{4}A_{n-4}
+2\fall{n}{3}A_{n-3}
-3\fall{n}{2}A_{n-2}
-2nA_{n-1}\right).
\]

At \(a=1\), \eqref{eq:stirling-product} identifies the first spectral
term with first-kind Stirling numbers. These are also the cycle-count
coefficients of permutations. The reciprocal-gamma factor is thus a
classical mechanism, closely related to mod-Poisson asymptotics
\cite{KowalskiNikeghbali}. The contribution to this programme is the
explicit spectral transfer, the uniform second coefficient estimate,
and the exact rational enclosure, not a claim that the
reciprocal-gamma phenomenon itself is new.
